\section{Idea 2: runs are positions minus edges}\label{sec:fibres}
\begin{definition}[Insertion fibres, gaps and swaps]\label{def:gaps}
For \(P\in D_{-x}\), its insertion fibre \(F(x,P)\) is the set of
slots \(0,\ldots,n-1\) in which reinserting \(x\) gives a member of
\(D\). Split the occupied slots into consecutive runs. A fibre with
\(r\) runs contributes \(r-1\) to \(G(D)\). Thus a gap counts a
separation between runs, rather than every missing slot.
Let \(A(D)\) count pairs \((\sigma,j)\) for which interchanging the
adjacent alternatives in positions \(j,j+1\) stays in \(D\).
An undirected swap edge contributes twice to \(A\).
\end{definition}
\lean{\path{Condorcet.insertionSlots}, \path{Condorcet.insertionWord},
\path{SlotWord.runs}, \path{SlotWord.gaps}, \path{Condorcet.gapCount},
\path{directedEdges}, \path{adjacencyCount}. Some local definitions
use \(\mathrm{Option}(X)\): \texttt{none} denotes the inserted label
and \texttt{some} the surviving labels. The transport back to labelled
orders is proved, rather than assumed.}

\begin{theorem}[Exact fibre identity]\label{thm:fibre}
For every finite order family, without any Condorcet hypothesis,
\[
 \boxed{\Sig(D)+G(D)+A(D)=nM.}
\]
\end{theorem}
\begin{proof}
A fibre with \(m\) occupied slots and \(r\) runs has \(m-r\) pairs
of consecutive occupied slots: each run of size \(s\) has \(s-1\)
edges. Summing occupied slots over all fibres counts every order once
for every deleted label, giving \(nM\). Summing runs gives \(\Sig+G\),
since every nonempty fibre contributes one output and \(r-1\) gaps.
Each consecutive-slot edge moves \(x\) past its neighbour. An actual
swap of adjacent \(a,b\) appears once with \(a\) distinguished and
once with \(b\) distinguished, precisely the twice-per-edge convention
of \(A\). Hence the sum of \(m-r\) is \(A\), proving the identity.
\end{proof}
\lean{\path{Condorcet.deletion_gap_adjacency_identity} states
\(\texttt{deletionSum}+\texttt{gapCount}+\texttt{adjacencyCount}
=(k+1)|D|\). The equality between occupied-slot edges and actual
swaps is proved by \path{slotAdjacencyCount_eq_adjacencyCount};
the local run identity is \path{SlotWord.nonempty_run_identity}.}

In Example~\ref{ex:four}, the swap graph is the path
\(123\leftrightarrow213\leftrightarrow231\leftrightarrow321\).
Thus \(A=6,G=0\), and \(6+0+6=3\cdot4\).
By contrast, \(\{123,231\}\) has \(\Sig=5,G=1,A=0\).
It is a peak-pit subset, but its 1-fibre above 23 has slots \(\{0,2\}\).
Completion matters in the next lemma.

\begin{lemma}[Interval fibres]\label{lem:interval}
For a full extreme-law-plus-pair family, every nonempty insertion fibre
is an interval. Therefore \(G=0\), and \(\Sig=nM-A\).
\end{lemma}
\begin{proof}
Fix the relative order \(P\) of all labels except \(x\). Every defining
constraint permits an interval of positions for \(x\). For instance,
\(xN_1\) forces \(x\) below the higher of the other two labels;
\(xN_3\) forces it above the lower. If a different label \(y\) is
never first, either the third label is already above \(y\), making
the condition automatic, or \(x\) must lie above \(y\). Never-last
is the reverse case. A fixed comparison involving \(x\) permits a
prefix or suffix of slots; constraints not involving \(x\) are constant.
Intersect these intervals. Fullness says that every surviving slot is
present, so the fibre has at most one run. Apply Theorem~\ref{thm:fibre}.
\end{proof}
\lean{\path{Condorcet.IsExtremeFull.insertionWord_interval},
\path{IsExtremeFull.gapCount_zero} and
\path{IsExtremeFull.deletion_plus_slotAdjacency}. The premise is
\path{IsExtremeFull}, not merely \path{IsPeakPit}.}
